% from ega4-2-fr.tex
\oldpage[IV-2]{125}

croissante de parties fermées $T$ de $Z$, de réunion $Z$),
$i_*(\mathcal F\vert X-T)$ (où $i:X-T\to X$ est l'injection canonique) est
cohérent.
\end{enumerate}
Lorsque $X$ vérifie l'une des hypothèses (i), (ii) de (5.11.1), a) et d) sont
équivalentes.
\end{corollary}

Notons que l'on a
$\operatorname{Ass}(i_*(\mathcal F\vert X-T))=\operatorname{Ass}(\mathcal F)$
en vertu de l'hypothèse et de (3.1.13); il résulte donc de (5.10.2) que les
applications canoniques
\[
i_*(\mathcal F\vert X-T)\longrightarrow\mathcal H^0_{X/Z}(\mathcal F)
\]
sont injectives; le fait que a) implique d) est donc conséquence de cette
remarque. Inversement, la condition d) implique, en vertu de (5.11.1), que
$\operatorname{codim}(T\cap Y_\alpha,Y_\alpha)\geq2$ avec les notations de
(5.11.1); par suite on a
$\operatorname{codim}(Z\cap Y_\alpha,Y_\alpha)\geq2$ puisque $Z$ est réunion
de ses parties qui sont fermées dans $X$, et la dernière assertion du
corollaire découle de (5.11.1).

\begin{corollary}[5.11.6]
\label{IV.2.s5.11.6-fr}
Soient $A$ un anneau noethérien, $X=\operatorname{Spec}(A)$. Considérons les
propriétés suivantes:
\begin{enumerate}
\item[a)] Pour tout anneau quotient intègre $B$ de $A$, l'anneau $B^{(1)}$
(notation de (5.11.2)) est une $B$-algèbre finie.
\item[b)] Pour tout $\mathcal O_X$-Module cohérent $\mathcal F$ et toute partie
$Z$ de $X$, stable par spécialisation, et telle que pour tout
$x\in\operatorname{Ass}(\mathcal F)\cap(X-Z)$ on ait
$\operatorname{codim}(\overline{\{x\}}\cap Z,\overline{\{x\}})\geq2$, le
$\mathcal O_X$-Module $\mathcal H^0_{X/Z}(\mathcal F)$ est cohérent.
\item[c)] Pour toute partie fermée $T$ de $X$ et tout $\mathcal O_U$-Module
cohérent $\mathcal G$ (où $U=X-T$) tel que pour tout
$x\in\operatorname{Ass}(\mathcal G)$ on ait
$\operatorname{codim}(\overline{\{x\}}\cap T,\overline{\{x\}})\geq2$,
$i_*(\mathcal G)$ (où $i:U\to X$ est l'injection canonique) est un
$\mathcal O_X$-Module cohérent.
\item[d)] Pour tout anneau quotient intègre $B$ de $A$ et tout idéal
$\mathfrak J$ de hauteur $\geq2$ dans $B$, l'anneau
\[
\bigcap_{\mathfrak p\not\supset\mathfrak J}B_\mathfrak p
\]
est une $B$-algèbre finie.
\end{enumerate}
On a alors les implications
\[
\mathrm a)\Longleftrightarrow\mathrm b)\Longrightarrow\mathrm c)
\Longleftrightarrow\mathrm d).
\]
En outre, les conditions a), b), c) et d) sont vérifiées dans chacun des deux
cas suivants:
\begin{enumerate}
\item[(i)] $A$ est quotient d'un anneau régulier.
\item[(ii)] $A$ est universellement caténaire et universellement japonais.
\end{enumerate}
\end{corollary}

Soit $B=A/\mathfrak q$, où $\mathfrak q$ est un idéal premier de $A$, de sorte
que $Y=\operatorname{Spec}(B)$ est la partie fermée $V(\mathfrak q)$ de $X$;
posons $Z=Z^{(2)}(Y)$, qui est une partie de $X$ stable par spécialisation;
comme $B$ est intègre, $\operatorname{Ass}(A/\mathfrak q)$ est réduit au point
générique $\mathfrak q$ de $Y$. Si la condition b) est vérifiée, on peut
l'appliquer au $\mathcal O_X$-Module cohérent
$\mathcal F=(A/\mathfrak q)^\sim$ et à $Z$, et en vertu de (5.9.3.1), cela
montre que $B^{(1)}$ est un $A$-module de type fini, et a fortiori un
$B$-module de type fini. Inversement, supposons a) vérifiée; alors, si
$\mathcal F$ est un $\mathcal O_X$-Module cohérent tel que pour tout
$x\in\operatorname{Ass}(\mathcal F)\cap(X-Z)$ on ait
$\operatorname{codim}(\overline{\{x\}}\cap Z,\overline{\{x\}})\geq2$, on peut
appliquer (avec les notations de (5.11.1)) à chacun des schémas affines
$Y_\alpha=\operatorname{Spec}(B_\alpha)$, où $B_\alpha$ est un anneau quotient
intègre de $A$, le résultat de a); comme par hypothèse $Z_\alpha$ est contenu
dans $Z^{(2)}(Y_\alpha)$, la condition a) (compte tenu de (5.10.2)

